\chapter{Blindness}

\section*{Definition}
Blindness and vision loss range from mild impairment to complete absence of light perception. Definitions and legal thresholds vary by country; clinicians describe vision using measured acuity, visual fields, and the effect on daily activities rather than a single universal cutoff.

\section*{When to See a Doctor}

\begin{itemize}
 \item Any sudden loss or marked change in vision, whether painful or painless, requires emergency assessment because retinal, vascular, optic-nerve, or neurological causes may be time-sensitive.
 \item Acute vision loss with eye pain, redness, or headache (suspect acute angle-closure glaucoma or optic neuritis)
 \item Transient vision loss (amaurosis fugax) concerning for embolic disease
 \item Progressive bilateral visual decline affecting daily function
\end{itemize}

\section*{How It Develops}
Blindness results from disruption of the visual pathway at any point: the transparent media (cornea, lens, vitreous), retina, optic nerve, optic chiasm, optic radiations, or occipital cortex. The mechanism depends on the causes and may involve neuronal apoptosis, ischemic infarction, demyelination, or mechanical compression.

\section*{Causes}
\begin{itemize}
 \item \textbf{Cataract:} Age-related lens opacification causing progressive light scatter
 \item \textbf{Glaucoma:} Chronic open-angle leading to irreversible retinal ganglion cell loss
 \item \textbf{Diabetic retinopathy:} Proliferative disease with vitreous hemorrhage and tractional detachment
 \item \textbf{Macular degeneration:} Neovascular (wet) or atrophic (dry) photoreceptor destruction
 \item \textbf{Trauma:} Open-globe injury, traumatic optic neuropathy, or hyphema
 \item \textbf{Infection:} Keratitis, endophthalmitis, or cytomegalovirus retinitis (immunocompromised)
 \item \textbf{Vascular:} Central retinal artery occlusion, central retinal vein occlusion, or anterior ischemic optic neuropathy
 \item \textbf{Neurologic:} Stroke affecting the occipital lobe or optic radiations
\end{itemize}

\section*{Related Symptoms}
\begin{itemize}
 \item Scotoma (central or peripheral)
 \item Photopsia
 \item Metamorphopsia (distortion of straight lines)
 \item Nyctalopia (night blindness)
 \item Dyschromatopsia (color vision deficiency)
\end{itemize}
\textbf{Important:} The World Health Organization estimates that at least 1 billion cases of vision impairment could have been prevented or remain unaddressed. Uncorrected refractive error and cataract are leading causes worldwide.

\section*{Medical Evaluation}
\begin{itemize}
 \item Measurement of best-corrected visual acuity and confrontation visual fields
 \item Pupillary examination (relative afferent pupillary defect indicates optic nerve pathology)
 \item Slit-lamp biomicroscopy and dilated funduscopy
 \item Optical coherence tomography (OCT) for macular and optic nerve assessment
 \item Humphrey or Goldmann perimetry for formal visual field mapping
 \item Neuroimaging (CT or MRI) if intracranial lesion or stroke suspected
\end{itemize}

\section*{Treatment Options}
\begin{itemize}
 \item Cataract extraction with intraocular lens implantation for lens opacification
 \item Intraocular pressure reduction (medications, laser, or surgery) for glaucoma
 \item Anti-VEGF intravitreal injections for neovascular macular degeneration and diabetic retinopathy
 \item Panretinal photocoagulation for proliferative diabetic retinopathy
 \item Acute retinal artery occlusion requires immediate ophthalmic and stroke-centre assessment; treatments such as ocular massage or thrombolysis are not routine and depend on specialist protocols and timing.
 \item Urgent corticosteroids are used when giant-cell arteritis is suspected, while optic neuritis and other inflammatory causes require specialist-directed treatment.
 \item Low-vision rehabilitation and assistive technology for permanent vision loss
\end{itemize}

\section*{Self-Care and Home Management}
\begin{itemize}
 \item Optimize home lighting and reduce glare for residual vision
 \item Use of low-vision aids (magnifiers, telescopic lenses, large-print materials)
 \item Contrast-enhancing strategies, such as dark cutting boards for light foods and tactile markings on appliances
 \item Training in orientation and mobility (cane use, guide dog)
 \item Social work referral for blind and low-vision support services
\end{itemize}

\section*{References}
\begin{thebibliography}{99}
\bibitem{blindness_vision_loss:ref1} Flaxman SR, Bourne RRA, et al. ``Global causes of blindness and distance vision impairment.'' \textit{Lancet Glob Health}. 2017;5(12):e1221-e1234.
\bibitem{blindness_vision_loss:ref2} Gilmore MS, Aflaki M, et al. ``Cataract.'' \textit{Lancet}. 2023;401(10383):1195-1208.
\bibitem{blindness_vision_loss:ref3} Quigley HA. ``Glaucoma.'' \textit{Lancet}. 2011;377(9774):1367-1377.
\bibitem{blindness_vision_loss:ref4} Kanski JJ, Bowling B. \textit{Kanski\'s Clinical Ophthalmology}, 9th ed. Elsevier, 2020.
\bibitem{blindness_vision_loss:ref5} World Health Organization. ``Blindness and vision impairment.'' Updated February 2026. https://www.who.int/news-room/fact-sheets/detail/blindness-and-visual-impairment.
\bibitem{blindness_vision_loss:ref6} American Academy of Ophthalmology. ``Central Retinal Artery Occlusion.'' EyeWiki, accessed September 2026. https://eyewiki.aao.org/Central\_Retinal\_Artery\_Occlusion.
\end{thebibliography}
